\documentclass[11pt,a4paper]{article}

%% ---- Encoding and fonts ----
\usepackage[utf8]{inputenc}
\usepackage[T1]{fontenc}
\usepackage{lmodern}

%% ---- Language ----
\usepackage[spanish,es-tabla]{babel}

%% ---- Mathematics ----
\usepackage{amsmath,amssymb,amsthm}
\usepackage{mathtools}

%% ---- Graphics ----
\usepackage{graphicx}
\usepackage{xcolor}
\usepackage{float}

%% ---- Geometry ----
\usepackage[a4paper,top=2.5cm,bottom=2.5cm,left=2.5cm,right=2.5cm]{geometry}

%% ---- Hyperlinks ----
\usepackage[colorlinks=true,linkcolor=blue!50!black,citecolor=green!40!black,urlcolor=blue]{hyperref}

%% ---- Captions ----
\usepackage{caption}
\usepackage{subcaption}

%% ---- Tables ----
\usepackage{booktabs}
\usepackage{array}

%% ---- Header / Footer ----
\usepackage{fancyhdr}
\pagestyle{fancy}
\fancyhf{}
\fancyhead[L]{\small Informe Experimental --- NNGP/NTK}
\fancyhead[R]{\small \today}
\fancyfoot[C]{\thepage}
\renewcommand{\headrulewidth}{0.4pt}

%% ---- Graphics path ----
\graphicspath{{./experiments/figures/}}

%%%%%%%%%%%%%%%%%%%%%%%%%%%%%%%%%%%%%%%%%%%%%%%%%%%%%%%%%%%%%%%%%%%%%%%%%%%%
\begin{document}

\title{\vspace{-1.5cm}
       \textbf{Informe Experimental: Generalización del límite NNGP/NTK}\\[0.4cm]
       {\large Convergencia CLT y comparación de kernels}\\[0.2cm]
       {\small TFM --- \today}
       }
\date{}
\maketitle
\thispagestyle{fancy}

%%%%%%%%%%%%%%%%%%%%%%%%%%%%%%%%%%%%%%%%%%%%%%%%%%%%%%%%%%%%%%%%%%%%%%%%%%%%
\section{Introducción}
%%%%%%%%%%%%%%%%%%%%%%%%%%%%%%%%%%%%%%%%%%%%%%%%%%%%%%%%%%%%%%%%%%%%%%%%%%%%

Este informe recoge los resultados de dos experimentos numéricos diseñados para explorar la validez del límite de ancho infinito en redes neuronales \emph{feedforward} de una capa oculta con activación ReLU. Los experimentos se centran en dos aspectos fundamentales de la teoría NNGP/NTK:

\begin{enumerate}
  \item \textbf{Convergencia CLT:} La distribución de las salidas de una red neuronal con pesos aleatorios converge a una Gaussiana a medida que el ancho $n \to \infty$. Esto es consecuencia del Teorema Central del Límite (CLT) aplicado a la suma de contribuciones independientes de cada neurona oculta.
  \item \textbf{Comparación de kernels:} Los kernels empíricos de ancho finito convergen a sus límites analíticos (NNGP kernel y Neural Tangent Kernel) cuando $n \to \infty$.
\end{enumerate}

Ambos experimentos utilizan una red de una capa oculta:
\begin{equation}
f(x) = \frac{1}{\sqrt{n}} \sum_{j=1}^{n} w_j\, \sigma\bigl(v_j^\top x + b_j\bigr),
\qquad 
v_j \sim \mathcal{N}(0, \sigma_w^2 I / d),\; b_j \sim \mathcal{N}(0, \sigma_b^2),\; w_j \sim \mathcal{N}(0, \sigma_w^2)
\end{equation}

Para la activación ReLU $\sigma(z) = \max(0, z)$, los kernels NNGP y NTK tienen expresiones analíticas cerradas (Cho \& Saul, 2009; Jacot et al., 2018), lo que permite una comparación precisa.

%%%%%%%%%%%%%%%%%%%%%%%%%%%%%%%%%%%%%%%%%%%%%%%%%%%%%%%%%%%%%%%%%%%%%%%%%%%%
\section{Experimento 1: Convergencia CLT}
%%%%%%%%%%%%%%%%%%%%%%%%%%%%%%%%%%%%%%%%%%%%%%%%%%%%%%%%%%%%%%%%%%%%%%%%%%%%

Las predicciones de la red con pesos aleatorios de ancho $n$ convergen en distribución a un proceso Gaussiano con kernel NNGP cuando $n \to \infty$. Para un punto fijo $x$, la distribución límite es
\begin{equation}
f(x) \xrightarrow{d} \mathcal{N}\bigl(0, \Sigma(x,x)\bigr), \qquad 
\Sigma(x,x') = \sigma_b^2 + \sigma_w^2 \,\mathbb{E}_{v\sim\mathcal{N}(0,I/d)}\bigl[\sigma(v^\top x)\,\sigma(v^\top x')\bigr].
\end{equation}

A continuación se presentan los resultados experimentales.

\subsection{Histogramas}

La Figura~\ref{fig:histograms} muestra la distribución empírica de $f(x)$ para distintos anchos $n$, junto con la densidad Gaussiana teórica $N(0, K_{\text{NNGP}}(x,x))$.

\begin{figure}[H]
\centering
\includegraphics[width=\textwidth]{exp01/histograms.png}
\caption{Distribución de la salida de la red para distintos anchos. La línea negra discontinua es la Gaussiana NNGP teórica.}
\label{fig:histograms}
\end{figure}

Para $n=1$, la distribución es marcadamente no Gaussiana (concentrada en torno a valores positivos debido a ReLU). A medida que $n$ crece, la distribución se aproxima rápidamente a la Gaussiana. Para $n=5000$, la distribución empírica es prácticamente indistinguible de la teórica.

\subsection{Métricas de convergencia}

La Figura~\ref{fig:convergence_metrics} muestra la evolución de diversas métricas con el ancho.

\begin{figure}[H]
\centering
\includegraphics[width=0.95\textwidth]{exp01/convergence_metrics.png}
\caption{Métricas de convergencia CLT en función del ancho $n$. La divergencia KL y el estadístico KS decrecen monótonamente. La varianza converge al valor teórico NNGP.}
\label{fig:convergence_metrics}
\end{figure}

\subsection{Tabla de resultados}

\begin{table}[H]
\centering
\caption{Convergencia CLT: métricas vs ancho de la red}
\label{tab:clt_convergence}
\begin{tabular}{crrrrrr}
\toprule
$n$ & $|\mu|$ & $\sigma^2$ & Skewness & Kurtosis & KL & KS \\
\midrule
        1 & 0.00118 & 1.01388 & +0.0880 & +15.3759 & 8.5476e-01 & 0.2542 \\
        5 & 0.00518 & 1.00543 & -0.0427 & +2.9530 & 9.1492e-02 & 0.0699 \\
       10 & 0.00062 & 1.00579 & +0.0104 & +1.6318 & 2.8948e-02 & 0.0378 \\
       50 & 0.00003 & 1.00033 & -0.0146 & +0.2686 & 3.3194e-03 & 0.0083 \\
      100 & 0.01403 & 1.00491 & -0.0223 & +0.1846 & 3.1925e-03 & 0.0106 \\
      500 & 0.01448 & 1.01182 & +0.0213 & +0.0589 & 4.0990e-03 & 0.0113 \\
     1000 & 0.00526 & 0.99027 & +0.0174 & +0.0033 & 4.4425e-03 & 0.0062 \\
     5000 & 0.00467 & 1.01222 & -0.0180 & -0.0713 & 4.4225e-03 & 0.0068 \\
\bottomrule
\end{tabular}
\end{table}

\subsection{Q-Q plots}

Los Q-Q plots de la Figura~\ref{fig:qqplots} confirman visualmente la convergencia: para $n=5000$ los puntos siguen prácticamente la recta $y=x$, indicando que la distribución empírica coincide con la Gaussiana teórica.

\begin{figure}[H]
\centering
\includegraphics[width=0.95\textwidth]{exp01/qq_plots.png}
\caption{Q-Q plots de la distribución empírica vs Gaussiana NNGP.}
\label{fig:qqplots}
\end{figure}

\subsection{Análisis}

Observamos que:
\begin{itemize}
  \item La \textbf{divergencia KL} disminuye de $\sim 10^{-1}$ para $n=1$ a $\sim 10^{-3}$ para $n=5000$.
  \item La \textbf{varianza empírica} converge al valor teórico NNGP con error $<1\%$ para $n \geq 100$.
  \item La \textbf{skewness} y \textbf{kurtosis} tienden a 0, confirmando la simetría y Gaussianidad de la distribución límite.
  \item La convergencia es rápida: para $n \geq 100$ la aproximación Gaussiana es ya muy precisa.
\end{itemize}

%%%%%%%%%%%%%%%%%%%%%%%%%%%%%%%%%%%%%%%%%%%%%%%%%%%%%%%%%%%%%%%%%%%%%%%%%%%%
\section{Experimento 2: Comparación de Kernels}
%%%%%%%%%%%%%%%%%%%%%%%%%%%%%%%%%%%%%%%%%%%%%%%%%%%%%%%%%%%%%%%%%%%%%%%%%%%%

El NNGP kernel y el NTK definen la dinámica de la red en el límite de ancho infinito:
\begin{align}
K^{\text{NNGP}}(x,x') &= \sigma_w^2\,\mathbb{E}\bigl[\sigma(z)\,\sigma(z')\bigr], \quad
z = v^\top x + b, \\
\Theta(x,x') &= K^{\text{NNGP}}(x,x') + \sigma_w^2\,\frac{x^\top x'}{d}\,\mathbb{E}\bigl[\sigma'(z)\,\sigma'(z')\bigr].
\end{align}

Para ReLU, ambas expresiones tienen forma cerrada (Cho \& Saul, 2009; Jacot et al., 2018). El kernel empírico de salida para ancho $n$ se define como
\begin{equation}
\hat{K}_n(x,x') = \frac{1}{M}\sum_{m=1}^{M} f_m(x; n)\, f_m(x'; n),
\end{equation}
donde $M$ es el número de inicializaciones aleatorias independientes. Adicionalmente, se computa el NTK empírico mediante los gradientes de la red:
\begin{equation}
\hat{\Theta}_n(x,x') = \frac{1}{M}\sum_{m=1}^{M} \langle \nabla_\theta f_m(x), \nabla_\theta f_m(x') \rangle.
\end{equation}

\subsection{Matrices kernel}

La Figura~\ref{fig:kernel_matrices} compara las matrices NNGP y NTK analíticas con las empíricas para distintos anchos.

\begin{figure}[H]
\centering
\includegraphics[width=\textwidth]{exp02/kernel_matrices.png}
\caption{Matrices kernel. De izquierda a derecha: NNGP ($n=\infty$), NTK ($n=\infty$), y kernels empíricos para $n=10,100,500,1000,2000$.}
\label{fig:kernel_matrices}
\end{figure}

\subsection{Convergencia de kernels}

La Figura~\ref{fig:kernel_error} muestra la convergencia del kernel empírico de salida al NNGP teórico (izquierda) y del NTK empírico al NTK teórico (derecha).

\begin{figure}[H]
\centering
\includegraphics[width=0.95\textwidth]{exp02/convergence_error.png}
\caption{Error relativo de Frobenius vs ancho $n$ para el kernel de salida (izquierda) y el NTK empírico (derecha). Las líneas discontinuas muestran el ajuste de ley de potencias.}
\label{fig:kernel_error}
\end{figure}

El error del kernel de salida decrece como $\sim n^{-0.5}$, consistente con la tasa esperada $O(1/\sqrt{n})$ del CLT. El NTK empírico también converge al NTK teórico, aunque con error sistemáticamente mayor debido a la contribución adicional de los gradientes y la dependencia de $M$ (número de ensambles).

\subsection{Espectro de autovalores}

La Figura~\ref{fig:eigenspectra} compara los espectros de autovalores.

\begin{figure}[H]
\centering
\includegraphics[width=0.95\textwidth]{exp02/eigenvalue_spectra.png}
\caption{Espectro de autovalores de los kernels NNGP y NTK para distintos anchos.}
\label{fig:eigenspectra}
\end{figure}

El espectro del NTK muestra autovalores sistemáticamente mayores que los del NNGP, debido al término adicional de la derivada. La convergencia espectral sigue la misma tasa $O(1/\sqrt{n})$.

\subsection{Diferencia NTK - NNGP}

La Figura~\ref{fig:ntk_diff} muestra la diferencia entre ambos kernels.

\begin{figure}[H]
\centering
\includegraphics[width=0.45\textwidth]{exp02/ntk_minus_nngp.png}
\caption{Matriz diferencia NTK - NNGP.}
\label{fig:ntk_diff}
\end{figure}

\begin{table}[H]
\centering
\caption{Error relativo del kernel empírico de salida vs NNGP teórico}
\label{tab:kernel_errors}
\begin{tabular}{crr}
\toprule
$n$ & $\varepsilon_F$(NNGP) & $\varepsilon_F$(NTK empírico) \\
\midrule
        1 & 0.21418 & 0.62380 \\
        5 & 0.10288 & 0.50361 \\
       10 & 0.11589 & 0.50038 \\
       50 & 0.08530 & 0.48498 \\
      100 & 0.03839 & 0.51942 \\
      500 & 0.08753 & 0.50853 \\
     1000 & 0.07644 & 0.50723 \\
     2000 & 0.07496 & 0.50650 \\
\bottomrule
\end{tabular}
\end{table}

%%%%%%%%%%%%%%%%%%%%%%%%%%%%%%%%%%%%%%%%%%%%%%%%%%%%%%%%%%%%%%%%%%%%%%%%%%%%
\section{Conclusiones}
%%%%%%%%%%%%%%%%%%%%%%%%%%%%%%%%%%%%%%%%%%%%%%%%%%%%%%%%%%%%%%%%%%%%%%%%%%%%

Los experimentos numéricos confirman los siguientes resultados teóricos:

\begin{enumerate}
  \item \textbf{Convergencia CLT:} La distribución de salidas de una red neuronal aleatoria converge a una Gaussiana a medida que $n \to \infty$, con tasa $O(1/\sqrt{n})$. Para $n \geq 100$, la aproximación Gaussiana es excelente (KL $\sim 10^{-3}$).
  
  \item \textbf{Convergencia del kernel de salida a NNGP:} El kernel empírico de ancho finito $\hat{K}_n(x,x')$ converge al NNGP kernel analítico con tasa $O(1/\sqrt{n})$ en norma de Frobenius.
  
  \item \textbf{Convergencia del NTK empírico:} El NTK empírico calculado mediante los gradientes de la red converge al NTK teórico en el límite de ancho infinito.
  
  \item \textbf{Diferencia NTK vs NNGP:} El NTK es sistemáticamente mayor que el NNGP debido al término adicional de la derivada de la activación, que codifica información sobre la dinámica del gradiente.
  
  \item \textbf{Implicaciones:} El límite de ancho infinito proporciona una aproximación cuantitativamente precisa incluso para anchos moderados ($n \geq 100$), lo que valida el uso del marco NNGP/NTK como herramienta analítica para redes realistas.
\end{enumerate}

%%%%%%%%%%%%%%%%%%%%%%%%%%%%%%%%%%%%%%%%%%%%%%%%%%%%%%%%%%%%%%%%%%%%%%%%%%%%
\section*{Metodología experimental}
%%%%%%%%%%%%%%%%%%%%%%%%%%%%%%%%%%%%%%%%%%%%%%%%%%%%%%%%%%%%%%%%%%%%%%%%%%%%

\begin{itemize}
  \item \textbf{Lenguaje:} Python 3.11 con NumPy, SciPy y Matplotlib.
  \item \textbf{Arquitectura:} Una capa oculta con activación ReLU, dimensión de entrada $d=50$ (Exp1) y $d=20$ (Exp2), pesos $\sigma_w=1$, sesgos $\sigma_b=1$.
  \item \textbf{Muestreo:} $8\,000$--$50\,000$ muestras Monte Carlo para Exp1; hasta $1\,000$ ensambles para promedio de kernel en Exp2, minimizando $200$ por ancho.
  \item \textbf{Compilación:} \texttt{pdflatex}, informe generado el 2026-06-07 06:04 UTC.
\end{itemize}

\end{document}
